%% ma: L12-B13-0014
%% lop: 12
%% chuong: 4
%% bai: 13
%% dang: 12.13.2
%% loai: TN4
%% muc: TH
%% dap_an: "D"
%% nguon: "(NBV)-ĐỀ SỐ 5-MỖI NGÀY 1 ĐỀ THI 2025.pdf (D05, MỖI NGÀY 1 ĐỀ - Đề số 5), Phần I Câu 12"
%% kien_thuc: "Bài 13 (thể tích vật thể tính qua diện tích thiết diện); Bài 12 (tích phân)"
%% trang_thai: chinh-thuc
%% ghi_chu: "Trùng ý tưởng với dạng 12.13.2 nên nhập cùng dạng"
\begin{ex}
Cho vật thể $(T)$ giới hạn bởi hai mặt phẳng $x=1$ và $x=3$, biết rằng khi cắt vật thể bởi mặt phẳng vuông góc với trục $Ox$ tại điểm có hoành độ $x$ $(1\le x\le3)$ thì được thiết diện là một hình chữ nhật có độ dài hai cạnh là $3x$ và $\sqrt{3x^2-1}$. Thể tích của vật thể $(T)$ được tính theo công thức
\choice{$V=3\pi\displaystyle\int_1^3x\sqrt{3x^2-1}\,\mathrm dx$}{$V=\pi\displaystyle\int_1^3x\sqrt{3x^2-1}\,\mathrm dx$}{$V=\displaystyle\int_1^3x\sqrt{3x^2-1}\,\mathrm dx$}{\True $V=3\displaystyle\int_1^3x\sqrt{3x^2-1}\,\mathrm dx$}
\loigiai{Diện tích thiết diện là $S(x)=3x\cdot\sqrt{3x^2-1}$. Thể tích vật thể là $V=\displaystyle\int_1^3S(x)\,\mathrm dx=3\int_1^3x\sqrt{3x^2-1}\,\mathrm dx$ (không có thừa số $\pi$ vì thiết diện là hình chữ nhật).}
\end{ex}
